\centerline{\bf \Huge Признаки делимости II}

\begin{flushright}
        \scriptsize{52}
\end{flushright}

\problem{\textbf{1}} Делится ли число 32561698 на 12? Ответ объясните.

\problem{\textbf{2}} Запишите несколько раз подряд число 2013 так, чтобы получившееся число делилось на 9. Ответ объясните.

\problem{\textbf{3}} К числу 15 припишите слева и справа по одной цифре так, чтобы полученное число делилось на 15. Перечислите все возможные варианты.

\problem{\textbf{4}} Чтобы открыть сейф, нужно ввести код — семизначное число, состоящее из двоек и троек.
Сейф откроется, если двоек в коде больше, чем троек, а сам код делится и на 3, и на 4. Какой
код может открывать сейф?

\problem{\textbf{5}} У Джерри есть девять карточек с цифрами от 1 до 9. Он выкладывает их в ряд, образуя девятизначное число. Том выписывает на доску все 8 двузначных
чисел, образованных соседними цифрами (например, для числа 789456123 это числа 78, 89, 94,
45, 56, 61, 12, 23). За каждое двузначное число, делящееся на 9, Том отдает Джерри кусочек
сыра. Какое наибольшее количество кусочков сыра может получить Джерри?

\problem{\textbf{6}} Известно, что число 55*44** делится на 72. Какие цифры могут стоять вместо звёздочек? Перечислите все возможные варианты.

\problem{\textbf{7}} Найдите наименьшее натуральное число $N$, такое, что число 99$\times N$ состоит из одних троек.